\chapter{Implementare și contribuții personale}

Acest capitol descrie cum a fost transpusă în practică soluția prezentată în capitolele anterioare. Dacă în capitolul precedent am discutat cerințele și deciziile de proiectare, acum mă concentrez pe implementarea propriu-zisă a componentelor platformei WhereWeFishin. Scopul nu este prezentarea exhaustivă a codului sursă, ci evidențierea modulelor dezvoltate, alegerilor tehnice și contribuțiilor personale care au făcut posibilă integrarea unui sistem multi-rol, geospațial și asistat de inteligență artificială.

Implementarea este prezentată pe straturi tehnice: mai întâi arhitectura generală, apoi backend-ul, frontend-ul și serviciul de analiză media. În fiecare secțiune am evidențiat explicit deciziile care reflectă contribuția personală.

\section{Arhitectura implementată a soluției}

Implementarea urmează o structură distribuită pe componente cu responsabilități bine delimitate. Interfața este o aplicație Angular, backend-ul este un API HTTP în ASP.NET Core, persistența este gestionată prin Entity Framework Core, iar analiza media este externalizată unui serviciu Python separat. Această separare permite evoluția independentă a modulelor, reduce cuplarea între zone cu cerințe tehnice diferite și menține claritatea codului.

Contribuția personală este vizibilă încă de la nivelul arhitecturii: am ales să separ operațiile tranzacționale uzuale de prelucrarea media costisitoare. Față de o abordare monolitică, comunicarea explicită dintre frontend, API și serviciul de analiză a simplificat atât dezvoltarea incrementală, cât și depanarea fluxurilor complexe.

\section{Implementarea backend-ului}

Backend-ul este nucleul logic al aplicației și zona în care sunt concentrate majoritatea regulilor de business. Implementarea sa se bazează pe controllere specializate, servicii reutilizabile și un strat de persistență bazat pe repository-uri. Această structură separă logica de transport HTTP de logica domeniului și de accesul la date, reducând duplicarea și simplificând testarea.

Configurația generală a aplicației definește autentificarea JWT, limitarea dimensiunii cererilor, compresia răspunsurilor, politicile CORS, rate limiting-ul pentru endpoint-urile sensibile și clienții HTTP necesari integrării cu serviciul Python. Această zonă arată că backend-ul nu tratează doar operații CRUD, ci gestionează și securitatea, performanța și interoperabilitatea.

\subsection{Autentificare, autorizare și controlul accesului}

Modulul de autentificare separă clar rolurile și protejează operațiile sensibile. Autentificarea se bazează pe tokenuri JWT, iar accesul la resurse este controlat prin claims și reguli de autorizare reutilizabile. Această soluție menține un model coerent pentru toate tipurile de utilizatori și reduce cuplarea dintre controllere și detaliile privind identitatea curentă.

Endpoint-urile de autentificare gestionează înregistrarea, autentificarea, verificarea tokenului și recuperarea accesului. Rata cererilor este limitată pentru a reduce riscul de abuz. Informațiile despre utilizator sunt extrase centralizat din claims, evitând repetarea logicii în fiecare controller.

Am introdus extensii de autorizare pentru resursele dependente de un loc de pescuit, astfel încât verificarea dreptului de administrare să poată fi refolosită în module diferite — pontoane, angajați, repopulări. Aceasta a permis extinderea aplicației cu noi module fără a duplica logica de acces.

Figura~\ref{fig:flux-autentificare} ilustrează fluxul de autentificare, de la cererea clientului până la validarea tokenului JWT.

\begin{figure}[ht]
\centering
\begin{tikzpicture}[
  node distance=1.1cm and 0.7cm,
  every node/.style={font=\small},
  step/.style={draw, rounded corners, align=center, minimum width=3cm, minimum height=0.95cm, fill=gray!10},
  arrow/.style={-Latex, thick}
]
\node[step] (b1) {Client trimite\\credențiale};
\node[step, right=of b1] (b2) {Backend verifică\\username și parolă};
\node[step, right=of b2] (b3) {Generare\\token JWT};
\node[step, right=of b3] (b4) {Token trimis\\clientului};
\node[step, right=of b4] (b5) {Cereri ulterioare\\cu token în header};

\draw[arrow] (b1) -- (b2);
\draw[arrow] (b2) -- (b3);
\draw[arrow] (b3) -- (b4);
\draw[arrow] (b4) -- (b5);
\end{tikzpicture}
\caption{Fluxul de autentificare bazat pe tokenuri JWT}
\label{fig:flux-autentificare}
\end{figure}

\subsection{Locuri de pescuit, pontoane și organizarea resurselor}

Implementarea locurilor de pescuit este unul dintre modulele centrale ale backend-ului, deoarece aproape toate fluxurile importante pornesc din această entitate. Controllerul gestionează listarea, filtrarea și administrarea locațiilor, cu interogări diferențiate în funcție de rolul utilizatorului. Am mutat filtrarea autorizată cât mai aproape de nivelul interogării, astfel încât datele să nu fie încărcate integral și filtrate doar la nivel de răspuns.

Pentru pontoane am implementat un modul separat care asociază zone distincte unui loc de pescuit și menține o legătură directă între componenta geospațială din interfață și regulile de disponibilitate din backend. Pontonul devine astfel o resursă rezervabilă explicită, nu doar o informație descriptivă. Această abordare este importantă pentru scenarii reale în care o locație conține mai multe spații cu disponibilitate independentă.

Contribuția personală constă în modelarea coerentă a relației dintre entitatea principală și subzonele rezervabile, precum și în reutilizarea acelorași reguli de autorizare pentru toate operațiile de administrare.

\subsection{Workflow-ul de aplicare pentru manager}

Unul dintre cele mai interesante module pe care le-am implementat este fluxul de aplicare pentru rolul de manager. În locul creării directe a unui loc de pescuit, aplicația introduce un proces controlat: utilizatorul transmite datele locației propuse, iar administratorul decide aprobarea sau respingerea cererii.

Am definit explicit stările aplicației, am păstrat istoricul de evaluare și am asociat clar decizia administrativă de crearea resursei în sistem. Controllerul gestionează crearea cererii, actualizarea ei în condiții controlate, retransmiterea după respingere și aprobarea finală. Am limitat tranzițiile de stare pentru a evita inconsistențe: nu se poate edita o cerere deja aprobată și nu pot exista două cereri active pentru același utilizator.

Entitatea aplicației păstrează informații despre administratorul evaluator, momentul evaluării și motivul respingerii. Această transparență face fluxul ușor de urmărit.

Această zonă reprezintă o contribuție personală importantă deoarece combină modelul de domeniu, regulile de business și administrarea controlată a creșterii platformei.

\subsection{Rezervări, sesiuni de pescuit și integrarea plății}

Fluxul de rezervare este cel mai dens din punct de vedere al logicii de business. La nivel de implementare, sistemul nu creează pur și simplu o înregistrare, ci validează locul de pescuit, verifică disponibilitatea pentru intervalul ales, corelează opțional un ponton, calculează costul și tratează separat integrarea plății.

Legătura cu Stripe este separată clar de persistența sesiunii finale. Am generat un \emph{PaymentIntent} care păstrează metadate relevante despre utilizator, locație și interval. Această separare oferă o bază mai sigură pentru corelarea plății cu sesiunea de pescuit și pentru evitarea erorilor de sincronizare. Modelul intern al sesiunii normalizează temporal datele și păstrează o stare explicită — \emph{Pending}, \emph{Confirmed} sau \emph{Cancelled}.

Contribuția personală în acest modul este integrarea coerentă a regulilor de disponibilitate, a modelului de sesiune și a unei plăți externe, fără a pierde consistența internă a datelor.

\subsection{Angajați, recenzii și servicii auxiliare}

Backend-ul include și componente care completează partea operațională. Modulul de angajați permite asocierea utilizatorilor la locuri de pescuit, creând o delegare controlată a operațiilor fără ridicarea privilegiilor la nivel de manager. Sistemul de recenzii introduce feedback din partea utilizatorilor și poate beneficia de caching al răspunsurilor și de invalidare explicită atunci când datele sunt modificate.

Am ales ca serviciile auxiliare să nu blocheze fluxurile critice. Dacă trimiterea unui email eșuează, operația principală nu este anulată. Contribuția personală constă în transformarea unor funcții aparent secundare într-o infrastructură stabilă care îmbunătățește experiența fără a compromite funcționarea sistemului.

\section{Implementarea frontend-ului}

Frontend-ul WhereWeFishin este o aplicație Angular modernă, bazată pe componente standalone, routing cu încărcare lazy și guard-uri funcționale. Această alegere a permis organizarea interfeței în zone funcționale clare și separarea fluxurilor destinate utilizatorului obișnuit, managerului, angajatului și administratorului.

\subsection{Structura aplicației SPA și controlul accesului în interfață}

La baza implementării se află structura traseelor și guard-urile funcționale. Fiecare zonă importantă se încarcă lazy și este protejată în funcție de rolul utilizatorului. Configurarea interceptorilor HTTP permite atașarea automată a tokenului JWT și tratarea controlată a erorilor recurente, precum expirarea autentificării.

Contribuția personală constă în distribuirea logicii de acces între rutare, guard-uri și interceptoare, astfel încât comportamentul interfeței să rămână coerent cu regulile de securitate ale API-ului.

\subsection{Harta interactivă și explorarea locurilor de pescuit}

Una dintre cele mai importante componente frontend este pagina de explorare, construită în jurul hărții interactive. Am combinat date despre locații, geolocalizarea utilizatorului, rutare pe hartă și informații despre specii. Am dezvoltat helperi dedicați pentru markere personalizate, stilizarea elementelor geografice și construirea ferestrelor informative într-un mod reutilizabil.

Harta funcționează ca punct de intrare în platformă: utilizatorul identifică rapid locațiile disponibile, vede proximitatea față de poziția sa și decide ce loc merită explorat mai departe. Contribuția personală constă în integrarea componentei Leaflet cu datele reale ale aplicației și în transformarea hărții dintr-un element vizual într-un instrument central de navigare.

\subsection{Coșul, calendarul și integrarea Stripe în interfață}

Fluxul de rezervare are o componentă frontend complexă. Utilizatorul selectează intervale, vede zilele ocupate, combină mai multe elemente în coș și finalizează plata. Am implementat un calendar care evidențiază perioadele rezervate, tratează intervale parțiale și blochează selecțiile invalide. Integrarea Stripe încarcă dinamic elementele de plată și menține separată logica de confirmare a tranzacției de logica afișării.

Contribuția personală constă în rezolvarea unor probleme practice de UX: vizualizarea clară a indisponibilităților, evitarea selecțiilor greșite și corelarea coerentă dintre formularul de rezervare și răspunsurile backend-ului.

\subsection{Wizard-ul de aplicare pentru manager și dashboard-ul operațional}

Zona de manager include două componente cu valoare tehnică ridicată. Wizard-ul de aplicare este organizat pe mai mulți pași și combină date descriptive, alegerea locației pe hartă și o revizuire finală înainte de trimitere. Dashboard-ul managerului reunește administrarea pontoanelor, gestionarea angajaților, actualizarea datelor locației și vizualizarea informațiilor operaționale.

Dashboard-ul este relevant deoarece concentrează mai multe subfluxuri într-o interfață unitară. Se vede direct legătura dintre modelul de date, regulile de business și reprezentarea lor în UI. Contribuția personală constă în construirea unei interfețe capabile să deservească activități operaționale reale, nu doar scenarii demonstrative.

\subsection{Verificarea prin cod QR și funcțiile administrative}

Am implementat și un flux de verificare a sesiunilor prin cod QR. Componenta folosește camera dispozitivului pentru citirea codului și tratează datele rezervării într-o formă suficient de tolerantă pentru a gestiona variații ale structurii lor. Acest lucru a fost important pentru a permite un flux rapid și stabil pe dispozitive mobile.

Zona administrativă oferă funcții de supraveghere asupra utilizatorilor, locațiilor și aplicațiilor de manager. Integrarea tuturor tipurilor de fluxuri într-o singură aplicație SPA a necesitat o separare clară a permisiunilor și a traseelor, fără a pierde coerența navigării.

\section{Implementarea sistemului de analiză media}

Sistemul de analiză media este componenta care diferențiază cel mai clar WhereWeFishin de o aplicație web convențională de rezervări. Implementarea implică colaborarea dintre trei straturi: interfața Angular, backend-ul .NET care orchestrează fluxul și serviciul Python care execută analiza.

\subsection{Integrarea dintre backend și serviciul Python}

La nivel de backend, un serviciu dedicat trimite fișierele către aplicația Python, interpretează răspunsul și actualizează entitățile persistente. Controllerele expun endpoint-uri pentru încărcare, listare, detaliere și ștergere, cu reguli clare de autorizare.

Am introdus stări explicite pentru analiză: \emph{Pending}, \emph{Processing}, \emph{Completed} și \emph{Failed}. Acestea permit urmărirea ciclului de viață al procesării și oferă frontend-ului o bază clară pentru polling și afișarea rezultatelor. Contribuția personală constă în transformarea unui apel extern câtre un serviciu de viziune artificială într-un flux controlat, persistent și ușor de integrat.

\subsection{Pipeline-ul video: detecție, tracking și post-procesare}

În serviciul Python, analiza video pornește de la încărcarea modelului YOLO și configurarea tracker-ului ByteTrack. Fiecare cadru video este analizat pentru detectarea obiectelor relevante, iar identitatea peștilor este urmărită între cadre pentru a evita contorizarea multiplă a aceluiași exemplar. Traiectoriile vizuale ale obiectelor sunt păstrate, iar rezultatul final este reîncodat pentru a putea fi servit eficient în interfață.

Contribuția personală se reflectă în configurarea pragurilor de încredere, alegerea mecanismului de trail și legătura dintre rezultatele brute ale modelului și forma în care sunt stocate și afișate ulterior.

\subsection{Afișarea rezultatelor în frontend și tratarea stării asincrone}

În frontend, modulele pentru recunoaștere video și clasificare pe imagini gestionează nu doar upload-ul fișierelor, ci și o stare asincronă care evoluează în timp. Am implementat un mecanism de polling pentru verificarea rezultatelor, iar interfața tratează separat stările de așteptare, succes și eroare. Pentru zona video am introdus și optimizări: încărcare întârziată a elementelor media și paginarea listelor de rezultate.

Contribuția personală constă în conectarea consecventă a stărilor persistente din backend cu mecanismele de feedback din interfață, construind un flux complet care gestionează sănătatea serviciului, așteptarea rezultatului și prezentarea sa utilă pentru utilizator.

\section{Sinteza capitolului}

Implementarea WhereWeFishin arată cum o soluție analizată la nivel de cerințe și proiectare poate fi transpusă într-un sistem coerent, multi-strat și orientat spre utilizare reală. Backend-ul concentrează regulile de business și integrarea serviciilor externe, frontend-ul construiește experiența utilizatorului și orchestrează fluxurile operaționale, iar serviciul de analiză media extinde aplicația cu o componentă de viziune artificială integrată controlat.

Contribuția personală se regăsește în toate aceste zone: în proiectarea și implementarea workflow-ului pentru manageri, în formalizarea rezervărilor și a plăților, în construirea unei interfețe geospațiale consistente și în integrarea unui pipeline AI capabil să proceseze atât imagini, cât și videoclipuri.
